\section{Related Work}
\label{sec:related_work}

A growing body of work shows that LLMs perform worse when instructions are revealed gradually rather than given upfront. \citet{laban2025llmslostmultiturnconversation} decompose fully-specified instructions into single-fact shards and deliver one per turn. Across 15 models and six tasks, they find a 39\% average performance drop. Most of this drop comes from a rise in unreliability across repeated runs, not from a loss in best-case skill. A control condition that delivers the same content in one reworded turn recovers 95.1\% of single-turn performance. This isolates the multi-turn structure itself as the cause. \citet{he2024multiifbenchmarkingllmsmultiturn} extend this setup to three-turn, multilingual dialogues. All 14 models tested degrade with each additional turn. o1-preview, for example, falls from 0.877 accuracy at turn one to 0.707 at turn three. The authors attribute much of this decline to \textit{instruction forgetting}, where a model abandons a constraint it had already satisfied earlier. \citet{jia2026battleanotherprobingllms} push further with EvolIF, which lets a conversation run until a simulated user's patience is exhausted, rather than stopping at a fixed turn count. GPT-5 proves the most resilient model tested, holding a 66.40\% robustness score, 5.59 points ahead of Gemini-3-Pro. Still, every model shares a steep drop near turns 4--5 and again near turns 11--12. The authors tie this to the difficulty of tracking many constraints and topic shifts at once. \citet{li2025structflowbenchstructuredflowbenchmark} instead sort turn-to-turn relationships into six categories, including follow-up, refinement, and recall. Across 13 models, \textit{refinement}, where a new instruction must revise earlier content rather than simply extend it, is by far the hardest category. Even the strongest model, DeepSeek-v3, reaches only 0.80 accuracy on refinement, against 0.99 on follow-up. The authors explain this gap with a simple failure: models mix up refinement with follow-up. When a new turn asks them to drop an old constraint, they often keep it anyway. Together, these studies show that multi-turn degradation is real and common. But each one studies tasks with an objectively verifiable answer, such as word counts or format rules, checkable by script rather than judgment. Creative writing is left out of this line of work on purpose. The closest benchmark to ours is SEQUOR \citep{canaverde2026sequormultiturnbenchmarkrealistic}. It filters out any constraint judged subjective, and names ``write a creative response'' as its own example of what gets removed. Even so, SEQUOR finds that accuracy drops far more when constraints build up over a conversation than when they are all given at once: over 11\% for a single constraint sustained across many turns, versus more than 40\% when several constraints must be tracked at once. This gap motivates our twist-shard design directly. It suggests that the burden of juggling constraints, not just remembering facts, drives multi-turn difficulty. Our paper studies exactly the regime SEQUOR excludes: open-ended, subjective creative writing.

Evaluating open-ended creative text is harder than evaluating factual output. There is no single correct answer, and quality spans several dimensions. \citet{chhun2022humancriteriaautomaticmetrics} identify six human criteria for story quality: relevance, coherence, surprise, engagement, empathy, and complexity. We adopt five of these in our own protocol. A follow-up study, \citet{chhun-etal-2024-language}, tests whether LLMs can substitute for human annotators on these same criteria. They find that LLM judges outperform traditional automatic metrics at the system level, though the judges still struggle to justify their own ratings. \citet{atmakuru2024cs4measuringcreativitylarge} probe creativity from a different angle. Their CS4 benchmark varies the number of constraints in a single-turn story prompt, and measures the trade-off between constraint satisfaction and narrative coherence. Our sharding design can be read as the multi-turn version of their single-turn approach: we spread constraints across turns instead of stacking them in one prompt. \citet{fein2025litbenchbenchmarkdatasetreliable} introduce LitBench, a set of 2{,}480 human-labeled story comparisons drawn from Reddit. Even the strongest off-the-shelf judge they test, Claude 3.7 Sonnet, reaches only 73\% agreement with human preference. This falls well below what a purpose-trained reward model achieves. This result motivates our own human validation study alongside the LLM-judge protocol. Two recent benchmarks argue that no single rubric can capture creativity. CreativityPrism \citep{hou2026creativityprismcrossdomainevaluationframework} evaluates 17 models across eight tasks in three domains, scoring each on quality, novelty, and diversity. It finds that strong performance on one dimension rarely generalizes to another; novelty scores in particular correlate weakly, or not at all, with the other two. WritingBench \citep{wu2025writingbenchcomprehensivebenchmarkgenerative} takes the opposite approach. It generates instance-specific evaluation criteria for each writing task, rather than applying one fixed rubric. We follow neither strategy exactly. Our fixed, small criterion set trades away some domain-specific sensitivity, in exchange for a direct, controlled comparison between FULL and Sharded. Humor, one of our six domains, has its own evaluation literature. HumorRank \citep{ajayi2026humorranktournamentbasedleaderboardevaluating} ranks nine models through pairwise, theory-grounded judgments, rooted in the General Theory of Verbal Humor. It reports high cross-judge agreement (Kendall's $\tau = 0.889$). This supports our use of pairwise comparison as a validation measure, and reinforces the view of humor as dependent on mechanism and timing, not just content.

Even outside multi-turn settings, LLMs do not process long context uniformly. \citet{liu2023lostmiddlelanguagemodels} show that accuracy on retrieval-dependent tasks follows a U-shaped curve, based on where the relevant fact sits in the context. Performance is strong at the start and end, and sharply worse in the middle. This effect persists even on a pure key-value lookup task stripped of semantics, and holds across every model tested except Claude-1.3. While this was shown on static contexts, \citet{laban2025llmslostmultiturnconversation} observe a related vulnerability in multi-turn dialogues: models often anchor on assumptions made in early turns or drop constraints introduced along the way. Our Sharded condition restates every prior shard at every new turn. This design works against both position-based forgetting and multi-turn state decay, rather than adding to it.

We rely on an LLM judge rather than rule-based scoring, since story quality resists scripted verification. \citet{zheng2023judgingllmasajudgemtbenchchatbot} show that GPT-4, used as a pairwise judge, agrees with human preference 85\% of the time. This exceeds the 81\% agreement observed between two human judges. But they also find a clear position bias. Under the default prompt, GPT-4 returns the same verdict after swapping answer order only 65\% of the time; a weaker model, Claude-v1, does so only 23.8\% of the time. We address this directly, by randomizing the left-right order of responses in every pairwise judgment. \citet{shi-etal-2025-judging} confirm this bias at a much larger scale, testing 15 LLM judges across roughly 150{,}000 evaluation instances spanning 22 tasks. They show that position bias is not random: it varies systematically by judge and by task, and is influenced most strongly by how close in quality the two compared solutions are. This further justifies our randomized-order protocol. \citet{liusie-etal-2024-llm} show that pairwise comparison gives more consistent LLM judgments than absolute scoring on NLG evaluation. This is why we treat pairwise agreement as a complementary, comparative check, alongside pointwise scores as our primary, absolute measure. The same pattern holds for human annotators: \citet{kiritchenko-mohammad-2017-best} show that best-worst scaling gives more reliable ratings than Likert-style scales, at equal annotation cost. Judge reliability for long-form creative evaluation remains limited, even with current models. We report human validation results on a fixed sample, to verify that our judges' preferences agree with human judgment on identical material.
